% ============================================================
%  04_clinical_brief.tex
%  Autonomy, Capacity, and the Right to Exit
%  CLINICAL & MEDICAL ETHICS BRIEF
%  For clinical ethics committees, psychiatric professionals,
%  hospital administrators, and medical associations
%  Justin Bogner, 2026
%  Compile: xelatex 04_clinical_brief.tex (twice)
% ============================================================

\documentclass[11pt, letterpaper]{article}
\usepackage{campaign}
\usepackage{xspace}
\usepackage{colortbl}

\campaignheader%
  {Autonomy, Capacity, and the Right to Exit}%
  {Clinical Brief --- Bogner, 2026}

\begin{document}

\campaigntitle%
  {AUTONOMY, CAPACITY, AND THE RIGHT TO EXIT}%
  {A Clinical and Medical Ethics Brief on Capacity-Based MAiD}%
  {Clinical / Medical Ethics Brief}%
  {Justin Bogner}%
  {2026}

\suitenote

\begin{policybox}[title={\sffamily\bfseries\small\color{C@accent}
  Purpose and Audience}]
This brief is written for clinical ethics committees, psychiatrists, psychiatric
nurses, hospital administrators, and medical associations engaged with Medical Aid in
Dying policy. It addresses: (1)~the clinical coherence of a capacity-based eligibility
standard; (2)~the assessment protocol proposed under this framework; (3)~the empirical
record from existing capacity-based regimes; (4)~the implications for the therapeutic
relationship; (5)~professional liability considerations; and (6)~specific clinical
objections, including the depressive incapacity objection.
\end{policybox}

\bigskip

\section{The Clinical Incoherence of Current Frameworks}

\subsection{The Informed Refusal Doctrine and Its Inconsistent Application}

The foundational principle of modern medical ethics is informed consent: a competent
patient may refuse any treatment, including life-sustaining treatment, on the basis of
their understanding of the treatment and its outcomes, without prior utilization and
without justification to the treating clinician.\footnote{%
  This principle is codified in legislation across jurisdictions (\textit{Health Care
  Consent Act}, Ontario, 1996; \textit{Patient Self-Determination Act}, US, 1990) and
  affirmed in case law. See \textit{Malette v.\ Shulman} (1990), 72 O.R.\ (2d) 417
  (right to refuse emergency treatment); \textit{Re B} [2002] 2 All ER 449 (right to
  refuse ventilator).}

The clinical inconsistency in current MAiD frameworks is notable: this same doctrine is
not applied to the decision to seek self-directed death. A patient may refuse
chemotherapy without trial; they may not, in most jurisdictions, access MAiD without
demonstrating terminal prognosis or externally validated suffering. The asymmetry is
not clinically justified. It reflects historical and legislative contingency, not
principled medical ethics.

\subsection{Suffering as a Construct Lacking Clinical Operationalization}

``Unbearable suffering'' is a legislative construct that clinical staff are asked to
operationalize without the validated instruments and inter-rater reliability that
capacity assessment has achieved. The Dutch SCEN protocol and the Belgian regional
review commissions have produced working definitions, but these have not yielded
the kind of standardized assessment frameworks that exist for decisional capacity.
The result is that eligibility assessment under suffering-based frameworks is:
(a)~subject to inter-rater unreliability; (b)~sensitive to clinician bias toward or
against MAiD provision; (c)~dependent on the applicant's ability to articulate
subjective distress in terms intelligible to the evaluator.

Capacity, by contrast, is a clinical concept with an established empirical literature,
validated assessment instruments, and cross-disciplinary consensus.\footnote{%
  Appelbaum PS, Grisso T. Assessing patients' capacities to consent to treatment.
  \textit{N Engl J Med} 1988;319(25):1635--1638. Moye J, Marson DC. Assessment of
  decision-making capacity in older adults: an emerging area of practice and research.
  \textit{J Gerontol B Psychol Sci Soc Sci} 2007;62(1):P3--P11.}
Replacing suffering with capacity as the eligibility criterion does not lower clinical
standards; it \textit{raises} them by substituting a measurable, clinically grounded
concept for a legally constructed but clinically unmoored one.

\section{The Capacity Assessment Protocol}

\subsection{Operationalizing Capacity for MAiD Assessment}

The sovereignty model applies the four-element Appelbaum-Grisso framework to MAiD:

\begin{description}
  \item[Understanding] The applicant can accurately describe what death involves
    (permanent, irreversible cessation), what the proposed provision entails, and what
    the realistic outcomes of all presented alternatives are. Understanding is assessed
    through structured interview, not assumption.

  \item[Appreciation] The applicant understands how the general information applies to
    their specific situation. Appreciation does not require optimism about alternatives;
    it requires demonstrated recognition that the information is relevant to their case.
    The evaluating clinician distinguishes between a patient who says ``SSRIs don't work
    for anyone'' (failure of understanding) and one who says ``SSRIs have a documented
    response rate of approximately 50\% for moderate depression; my history includes two
    adequate trials without response'' (appreciation satisfied).

  \item[Reasoning] The applicant can articulate a coherent account of how their
    understanding of their situation, the alternatives, and their values leads to their
    expressed preference. The reasoning need not be philosophically sophisticated; it
    must be internally coherent and free from gross logical errors or fixed false beliefs.

  \item[Expression of a Consistent Choice] The preference is stable across the
    assessment period and across multiple evaluation contexts. Contextual variation
    (presenting differently to different evaluators) is diagnostically relevant.
    Persistent stability across 90 days and multiple independent evaluations satisfies
    this element.
\end{description}

\subsection{The Evaluating Clinician's Role}

The role of the evaluating clinician under this framework is fundamentally different
from the role assigned under suffering-based frameworks.

\begin{table}[h!]
\centering
\small
\renewcommand{\arraystretch}{1.5}
\begin{tabularx}{\textwidth}{>{\sffamily\bfseries\small}p{3.2cm} X X}
\toprule
\rowcolor{C@accent}
{\color{white}\sffamily\bfseries Dimension} &
{\color{white}\sffamily\bfseries Protectionist Model} &
{\color{white}\sffamily\bfseries Sovereignty Model} \\
\midrule
Primary question
  & Is this person suffering enough?
  & Does this person understand what they are deciding? \\
Role
  & Gatekeeper / adjudicator of suffering
  & Assessor of comprehension and stability \\
Treatment history
  & Required; exhaustion of alternatives typical prerequisite
  & Irrelevant to eligibility; comprehension of alternatives is assessed \\
Clinical instrument
  & None standardized; clinician judgment on suffering
  & Appelbaum-Grisso framework; standardized capacity assessment \\
Liability orientation
  & Institutional (avoid death on our watch)
  & Patient-centered (assess authentic preference) \\
Therapeutic role
  & Custodial retention
  & Restorative alliance \\
\bottomrule
\end{tabularx}
\caption{Comparison of clinician roles under Protectionist and Sovereignty models.}
\end{table}

\subsection{Why 90 Days}

The minimum 90-day assessment period is calibrated to clinical reality rather than
legislative convenience. Three considerations support it. First, an adequate
antidepressant trial requires 6--12 weeks; a stability assessment that runs shorter
than this cannot meaningfully evaluate whether the preference survives a fluctuation
cycle. Second, 90 days spans typical episodic variability in mood disorders, allowing
evaluators to observe the preference across periods of relative improvement and
relative worsening. Third, the 90 days is a floor, not a ceiling: complex presentations
warrant longer assessment, and the framework specifies no maximum. The Dutch and
Belgian regimes use ``durable wish'' language without numerical anchors; the proposal
here trades some flexibility for procedural clarity, on the judgment that a published
minimum is more defensible than an unspecified ``stability'' standard, which is
subject to evaluator variance.

\subsection{Treatment-Resistant Depression: A Moving Landscape}

The alternatives presentation requirement is dynamic, not static. The clinical landscape
for treatment-resistant depression has changed materially since the frameworks governing
most MAiD statutes were drafted. Evaluators must present the current state of the field,
not its 2016 equivalent. Response rate figures below are taken from the conservative
end of the available meta-analytic and pivotal-trial estimates.

\begin{description}
  \item[Esketamine (Spravato)] FDA-approved March 2019 for treatment-resistant
    depression. Intranasal administration; rapid onset (hours rather than weeks);
    response rates in TRD populations approximately 50--60\% in short-term trials, with
    statistically significant superiority over placebo plus active oral antidepressant
    (NNT for response approximately 6).\footnote{%
    Fedgchin M, Trivedi M, Daly EJ, et al. Efficacy and safety of fixed-dose esketamine
    nasal spray combined with a new oral antidepressant in treatment-resistant
    depression: results of a randomized, double-blind, active-controlled study
    (TRANSFORM-1). \textit{Int J Neuropsychopharmacol} 2019;22(10):616--630.
    Pooled meta-analytic estimates: response RR 1.28 (95\% CI 1.12--1.46), remission
    RR 1.39 (95\% CI 1.18--1.63).}
    Durability at 6--12 months is less established. REMS certification required for
    prescribers and dispensing facilities. An applicant who has not been presented
    with esketamine as an option has not received an adequate alternatives presentation
    for a TRD assessment.

  \item[Repetitive Transcranial Magnetic Stimulation (rTMS)] FDA-cleared for major
    depressive disorder; pooled response rates approximately 29\% in primary
    meta-analytic samples, with consistent statistical superiority over sham (NNT
    approximately 5).\footnote{%
    Berlim MT, van den Eynde F, Tovar-Perdomo S, Daskalakis ZJ. Response, remission and
    drop-out rates following high-frequency repetitive transcranial magnetic stimulation
    (rTMS) for treating major depression: a systematic review and meta-analysis of
    randomized, double-blind and sham-controlled trials. \textit{Psychol Med}
    2014;44(2):225--239.}
    More recent network meta-analyses confirm efficacy with somewhat higher
    point estimates depending on stimulation parameters. Non-invasive; no anesthesia
    required; outpatient protocol. Growing evidence base for accelerated TMS protocols
    (intensive daily sessions).

  \item[Electroconvulsive Therapy (ECT)] Remains the highest-response-rate intervention
    for severe depression: pooled response rates approximately 74\% in unipolar major
    depressive episode, with remission rates of approximately 52\%
    overall.\footnote{%
    Bahji A, Hawken ER, Sepehry AA, Cabrera CA, Vazquez G. ECT beyond unipolar major
    depression: systematic review and meta-analysis of electroconvulsive therapy in
    bipolar depression. \textit{Acta Psychiatr Scand} 2019;139(3):214--226.}
    Efficacy is reduced but still substantial in TRD specifically: remission rates of
    approximately 48\% in patients with prior pharmacotherapy failure compared with
    65\% in patients without.\footnote{%
    Heijnen WT, Birkenh{\"a}ger TK, Wierdsma AI, van den Broek WW. Antidepressant
    pharmacotherapy failure and response to subsequent electroconvulsive therapy:
    a meta-analysis. \textit{J Clin Psychopharmacol} 2010;30(5):616--619.}
    Underutilized due to stigma and access limitations, not efficacy. Durability
    requires continuation strategies; relapse without continuation pharmacotherapy or
    maintenance ECT is substantial. Evaluators presenting alternatives must include
    ECT with accurate efficacy and risk information, not a stigma-inflected summary.

  \item[Ketamine (IV)] Off-label; rapid antidepressant effect; response rates in TRD
    approximately 50--65\% in short-term trials. Access variable; mechanism distinct
    from conventional antidepressants; evidence base growing.
\end{description}

The framework's response to this landscape is precisely the one that avoids paternalism:
the question is not whether the applicant has tried these modalities, but whether they
understand what each offers and have declined following genuine comprehension. Informed
refusal of esketamine is as valid as informed refusal of a second SSRI. The evaluator's
obligation is to ensure the refusal is genuinely informed---which requires knowing that
esketamine exists.

\subsection{What the Clinical Interview Is Not}

The clinical interview under the sovereignty model is not:

\begin{itemize}
  \item A session designed to persuade the patient to choose differently. The clinician
    presents alternatives in sufficient specificity for genuine comprehension; they do
    not advocate for any alternative.
  \item An assessment of whether the clinician agrees with the patient's values or
    finds their reasoning compelling. Clinicians routinely work with patients whose
    values differ from their own without this constituting an eligibility barrier.
  \item An opportunity to require trial of interventions the patient has declined.
    Informed refusal of alternatives is binding. The clinical interview's role is
    to assess whether the refusal is informed, not to override it.
\end{itemize}

\section{The Empirical Record from Capacity-Based Regimes}

Any clinical brief proposing a capacity-based standard must engage with the empirical
record from the jurisdictions that already operate one. The relevant data come primarily
from the Netherlands and Belgium.

\subsection{The Kim 2016 Findings}

The most-cited empirical analysis of psychiatric MAiD outcomes is Kim, De Vries, and
Peteet's 2016 review of 66 Dutch psychiatric euthanasia cases from 2011 to
2014.\footnote{%
  Kim SYH, De Vries RG, Peteet JR. Euthanasia and assisted suicide of patients with
  psychiatric disorders in the Netherlands 2011 to 2014. \textit{JAMA Psychiatry}
  2016;73(4):362--368. doi:10.1001/jamapsychiatry.2015.2887.}
The findings are not flattering to the Dutch regime as practiced and must be addressed
directly.

The study found that 56\% of psychiatric EAS recipients were described in case reports
as socially isolated or lonely; 32\% had been refused EAS by at least one physician
before ultimately receiving it, with most of the redirected cases handled by a single
mobile End-of-Life Clinic; in 11\% of cases there was no independent psychiatric
consultation; and physician disagreement about eligibility was common but the regional
review committees generally deferred to the judgments of the providing physicians.
The companion analysis of capacity evaluations specifically (Doernberg, Peteet, and
Kim, 2016) reached similar conclusions about variability in how the four-element
framework was applied.\footnote{%
  Doernberg SN, Peteet JR, Kim SYH. Capacity evaluations of psychiatric patients
  requesting assisted death in the Netherlands. \textit{Psychosomatics} 2016;57(6):
  556--565. doi:10.1016/j.psym.2016.06.005.}

\subsection{What the Sovereignty Model's Protocol Is Designed to Address}

These findings are the procedural failure modes the protocol specified in this brief
is designed to prevent. The mapping is direct.

The Kim findings show that a substantial minority of cases were redirected to a single
willing-provider clinic after refusal elsewhere. The sovereignty model's response is
the multi-evaluator concurrence requirement (\S 4.3): a finding of capacity requires
reasoned concurrence on each functional element by at least two independent evaluators,
not a workaround through a willing third party.

The Kim findings document that 56\% of cases involved social isolation or loneliness as
a reported feature. The sovereignty model's response is the independent coercion
screening by an unaffiliated social worker or advocate, the financial and dependency
circumstances review, and the explicit socioeconomic documentation requirement where
inadequate support is identified as a driver. These are not optional. They are
required elements of every assessment.

The Kim findings document that 11\% of cases proceeded without independent psychiatric
consultation. The sovereignty model's response is that independent psychiatric
evaluation is mandatory, not advisory.

The Kim findings document that review committees deferred to physician judgment rather
than examining contested decisions de novo. The sovereignty model's response is the
graduated review structure: explicit documentation of any divergent findings, mandatory
Clinical Ethics Committee referral within fourteen days for unresolved disagreement,
and administrative review with documented reasoned findings if the disagreement
persists.

A clinical reader should not read the Dutch experience as evidence against
capacity-based assessment. The Dutch regime is not a controlled implementation of the
sovereignty model; it is a different regime, with documented procedural gaps, and it
shows precisely where a more rigorous protocol must close those gaps.

\section{The Depressive Incapacity Objection: Clinical Analysis}

This is the objection most frequently raised in clinical contexts, and it deserves
careful clinical analysis rather than brief dismissal.

\subsection{The Objection}

Persons with treatment-resistant depression cannot satisfy the ``appreciation'' criterion
for decisional capacity because hopelessness functions as a cognitive filter: alternatives
are formally understood but experienced as inaccessible. On this account, a capacity
assessment that returns ``capacity present'' in a person with severe depression is a
false positive---the capacity criteria have been formally satisfied, but the authentic
preconditions for autonomous preference are absent.

\subsection{The Clinical Response}

\subsubsection{Depression does not uniformly impair capacity}

The empirical literature does not support the proposition that depression, including
severe depression, uniformly impairs the functional elements of decisional capacity.
The MacArthur Treatment Competence Study found that significant impairments in
decisional abilities affected only a minority of patients with major depression, with
deficits less pronounced than among patients with schizophrenia.\footnote{%
  Grisso T, Appelbaum PS. The MacArthur Treatment Competence Study. III: Abilities of
  patients to consent to psychiatric and medical treatments. \textit{Law and Human
  Behavior} 1995;19(2):149--174.}
Vollmann et al.\ (2003) found that on physician clinical assessment 2.9\% of depressed
inpatients were judged incompetent, rising to approximately 20\% under formal MacCAT-T
administration with stricter cut-offs---meaning that even on the most demanding formal
standard, four out of five depressed inpatients retained capacity.\footnote{%
  Vollmann J, Bauer A, Danker-Hopfe H, Helmchen H. Competence of mentally ill
  patients: a comparative empirical study. \textit{Psychological Medicine} 2003;33(8):
  1463--1471.}
The 2013 systematic review by Hindmarch, Hotopf, and Owen reaffirmed that capacity
impairment in depression is variable, correlated with severity, and not reliably
predicted by diagnosis alone.\footnote{%
  Hindmarch T, Hotopf M, Owen GS. Depression and decision-making capacity for treatment
  or research: a systematic review. \textit{BMC Medical Ethics} 2013;14:54.
  doi:10.1186/1472-6939-14-54.}
Across studies, the consistent finding is that the majority of patients with
depression---including inpatient samples---retain decisional capacity for treatment
decisions on rigorous formal assessment.

The implication is direct: a categorical ``depression negates capacity'' inference is
not supported by the data. Capacity is decision-specific, varies by severity, and must
be assessed individually rather than presumed absent on diagnostic grounds.

\subsubsection{Appreciation does not require experiential access}

The appreciation criterion requires that the patient understand how the information
applies to their situation. It does not require that they \textit{experience} alternatives
as accessible. A patient who accurately states that SSRIs have a documented response
rate for their condition, that two adequate trials without response are in their record,
and that they decline further trials because the probability of benefit does not, in
their assessment, justify the cost---satisfies the appreciation criterion. The criterion
is epistemic, not phenomenological.

To require, additionally, that the patient \textit{feel hope} that alternatives might
succeed is to import a positive affective requirement that has no basis in the
Appelbaum-Grisso framework and no precedent in clinical capacity doctrine. It also
creates an impossible standard: hopelessness is a symptom of the condition being
assessed; requiring its absence as a precondition for capacity makes the capacity
determination circular.

\subsubsection{The double standard problem}

The depressive incapacity objection, applied consistently, would require that persons
with depression satisfy a higher capacity standard than persons without depression for
all medical decisions. But medicine does not currently impose this: a patient with
severe depression who accepts pharmacological treatment does not have their capacity
for that decision specially scrutinized. The heightened scrutiny applies only when the
decision is to refuse the state-preferred intervention---revealing that the objection
is not about capacity at all, but about the outcome.\footnote{%
  This asymmetric scrutiny is well-documented. See Astrachan IM, Keene AR, Kim SYH.
  Questioning our presumptions about the presumption of capacity. \textit{J Med Ethics}
  2024;50(7):471--475. doi:10.1136/jme-2023-109199. See also Owen GS, Szmukler G,
  Richardson G, et al.\ Mental capacity and psychiatric in-patients: implications for
  the new mental health law in England and Wales. \textit{Br J Psychiatry}
  2009;195(3):257--263.}

\subsection{The Clinical Protocol Response}

The sovereignty model does not ignore the depressive incapacity concern. It addresses
it procedurally. A preference that is generated by acute depressive crisis is unlikely
to be stable across 90 days of assessment. A preference that survives the assessment
period---including periods of relative improvement and relative worsening---and
that is articulated consistently across multiple independent evaluations, has survived
the procedural filter that the objection is concerned about. The process is the
response.

\section{Evaluator Disagreement: Clinical Governance}

The multi-evaluator requirement raises a governance question that must be addressed
explicitly: what happens when independent evaluators reach different conclusions about
the presence of decisional capacity?

\subsection{Why Disagreement Occurs}

Capacity assessments in complex psychiatric presentations are clinically challenging.
Reasonable evaluators applying the same criteria to the same patient may weigh the
functional elements differently: one assessor may find the applicant's appreciation
criterion satisfied while another finds evidence of cognitive distortion affecting
probability estimation. Pretending this does not occur is not clinical
honesty.\footnote{%
  Appel JM. Legal and ethics considerations in capacity evaluation for medical aid in
  dying. \textit{J Am Acad Psychiatry Law} 2024;52(3):311--326.
  doi:10.29158/JAAPL.240038-24.}

\subsection{The Protocol}

The proposed governance protocol for evaluator disagreement operates in three steps:

\begin{enumerate}
  \item \textbf{Documentation of disagreement.} Each evaluator produces a written
    assessment documenting their findings on each of the four functional elements.
    Disagreement is recorded explicitly, with reasoned findings on each element.
    A ``capacity present'' finding must be reasoned; an unexamined concurrence is not
    sufficient.

  \item \textbf{Clinical Ethics Committee referral.} Where evaluators reach divergent
    conclusions, the case is referred to an independent Clinical Ethics Committee (CEC)
    within fourteen days. The CEC's function is not to conduct a new capacity assessment
    but to review the divergent findings and identify whether the disagreement reflects
    a genuine clinical difference or a procedural gap. The CEC may recommend additional
    evaluation on a specific element (e.g., further assessment of the appreciation
    criterion using counter-evidence presentation).

  \item \textbf{Administrative review for persistent disagreement.} Where CEC review
    does not resolve the disagreement, the matter is referred to the administrative
    review panel. This is a legal and governance step, not a clinical one. It creates
    a documented record of the dispute and a reviewable administrative decision---
    which is accessible to the applicant via appeal.
\end{enumerate}

\subsection{What Capacity Requires}

A finding of capacity requires reasoned concurrence by at least two independent
evaluators on each of the four functional elements. Majority rule on a binary question
(``capacity present / absent'') is insufficient; reasoned concurrence on each element
is the standard. This is consistent with the clear and convincing evidence standard
governing the assessment and protects both the applicant and the clinical assessors.

\section{Liability Considerations}

\subsection{The Current Liability Model}

Current clinical practice around suicidal ideation is oriented primarily around
institutional liability management: documenting that intervention occurred and that the
clinician met the standard of care for suicide prevention. This orientation produces
defensive practice patterns documented in the existing literature: coerced safety
contracts of no demonstrated efficacy, pressured hospitalization optimized for
administrative rather than therapeutic outcomes, and discharge into unchanged conditions.

This orientation often prioritizes liability reduction over the therapeutic
alliance. This has consequences for the therapeutic relationship, for patient trust,
and for clinical outcomes.

\subsection{The Sovereignty Model and Liability Reform}

Under the sovereignty model, the clinical liability framework shifts. Clinicians who
conduct rigorous, documented capacity assessments following the protocol specified in
this framework---multiple independent evaluations, documented alternatives presentation,
independent coercion screening, 90-day assessment period---are operating within a
structured and documentable standard of care. Liability exposure is reduced by
procedural compliance, not by preventing death at all costs.

The conscientious objection provision is critical here: no clinician is required to
participate in MAiD assessment or provision. The only obligation imposed on objectors
is timely referral to a willing provider. This protects clinical autonomy while
ensuring that patient access is not foreclosed by institutional objection.

\subsection{The Therapeutic Consequence}

The most significant clinical consequence of the sovereignty model is the reorientation
of the therapeutic relationship. When the patient possesses an enforceable exit right,
the clinician loses the coercive power of custodial retention.

This constraint is not a limitation on clinical effectiveness; it is the precondition
for a different kind of clinical effectiveness. The therapeutic relationship becomes
genuinely collaborative. The clinician's role is to make the patient's life worth
continuing---to address the conditions that are driving the preference, to connect the
patient with resources they may not know exist, to work on the quality of the life
rather than the fact of its continuation. This is a more honest, and ultimately more
effective, clinical orientation.

\section{Assessment Protocol Summary}

\begin{findingbox}[title={\sffamily\bfseries\small\color{C@accent}
  Clinical Protocol at a Glance}]
\small
\textbf{Stage 1 --- Decisional Capacity Assessment} (multiple independent evaluators)
\begin{itemize}
  \item Understanding: permanence of death; realistic outcomes of all alternatives
  \item Appreciation: application of information to own situation
  \item Reasoning: coherent, internally consistent account
  \item Expression: stable preference across assessment period
  \item State assessment: no acute psychosis or crisis-state distortion \textit{currently}
\end{itemize}

\medskip
\textbf{Stage 2 --- Process Comprehension}
\begin{itemize}
  \item Applicant understands why safeguards exist and accepts them
  \item Demonstrates sustained engagement across 90-day minimum period
  \item Failure to engage with Stage 2 is diagnostically relevant to Stage 1
\end{itemize}

\medskip
\textbf{Required Elements of Every Assessment}
\begin{itemize}
  \item Documented alternatives presentation with referrals offered; declinations recorded
  \item Independent coercion screening by unaffiliated social worker or advocate
  \item Financial and dependency circumstances review
  \item Offer of independent legal assistance
  \item Socioeconomic documentation where inadequate support is identified as a driver
\end{itemize}

\medskip
\textbf{Revocation:} Absolute, unconditional, available until the moment of provision.
\end{findingbox}

\section{Clinical Objections: Summary Responses}

\begin{objectionbox}[title={\sffamily\bfseries\small\color{C@warning}
  ``We cannot assess capacity reliably in psychiatric populations''}]
\textbf{Response:} Capacity assessment in psychiatric populations is a well-established
clinical practice with validated instruments (MacArthur Competence Assessment Tool;
Aid to Capacity Evaluation). The empirical record demonstrates that the majority of
patients with depression retain capacity on rigorous formal assessment (Grisso \&
Appelbaum, 1995; Vollmann et al., 2003; Hindmarch et al., 2013). The proposed framework
applies these instruments to a new context; it does not require new clinical methodology.
The challenge of assessing capacity in complex cases is an argument for rigorous
multi-evaluator assessment, not for categorical exclusion of psychiatric populations
from eligibility.
\end{objectionbox}

\medskip

\begin{objectionbox}[title={\sffamily\bfseries\small\color{C@warning}
  ``Psychiatrists will face pressure to approve requests against their clinical judgment''}]
\textbf{Response:} The conscientious objection provision protects clinicians who
decline to participate. The multi-evaluator requirement means no single clinician
carries the assessment alone. The documented capacity protocol creates a professional
standard of care that protects clinicians who conduct rigorous assessments, regardless
of outcome. The pressure concern is addressed by process, not by maintaining a
categorical prohibition.
\end{objectionbox}

\medskip

\begin{objectionbox}[title={\sffamily\bfseries\small\color{C@warning}
  ``This will normalize suicide and undermine suicide prevention''}]
\textbf{Response:} Suicide prevention is directed at crisis-state, impulsive, and
untreated-pathology-driven deaths. The sovereignty model is addressed at a categorically
different population: persons with stable, long-considered, informed preferences who
have survived a 90-day assessment process. These populations do not overlap clinically.
Indeed, the explicit, documented nature of MAiD provision---as distinct from unassisted
suicide---may reduce rather than increase the broader social normalization concern,
since it is administered within a visible, regulated, documented clinical process.
\end{objectionbox}

\medskip

\begin{objectionbox}[title={\sffamily\bfseries\small\color{C@warning}
  ``The Dutch experience shows capacity-based regimes fail''}]
\textbf{Response:} The Dutch experience documented in Kim et al.\ (2016) illustrates
specific procedural failure modes---redirected cases through a single willing provider,
high rates of social isolation among recipients, occasional absence of independent
psychiatric evaluation, deferential review committees. The sovereignty model's protocol
addresses each of these directly: multi-evaluator concurrence on each functional
element, mandatory independent coercion and socioeconomic screening, mandatory
independent psychiatric evaluation, and graduated administrative review with reasoned
findings. The Dutch data are evidence for a more rigorous protocol, not against
capacity-based assessment as such.
\end{objectionbox}

\bigskip
{\small\sffamily\color{C@midgray}
\textbf{Key Clinical References:}\\
Appel JM (2024), \textit{J Am Acad Psychiatry Law} 52(3):311--326;\\
Appelbaum PS \& Grisso T (1988), \textit{N Engl J Med} 319(25):1635--1638;\\
Astrachan IM, Keene AR, Kim SYH (2024), \textit{J Med Ethics} 50(7):471--475;\\
Bahji A et al.\ (2019), \textit{Acta Psychiatr Scand} 139(3):214--226;\\
Beauchamp TL \& Childress JF, \textit{Principles of Biomedical Ethics} (8th ed., 2019);\\
Berlim MT et al.\ (2014), \textit{Psychol Med} 44(2):225--239;\\
Doernberg SN, Peteet JR, Kim SYH (2016), \textit{Psychosomatics} 57(6):556--565;\\
Fedgchin M et al.\ (2019), \textit{Int J Neuropsychopharmacol} 22(10):616--630;\\
Grisso T \& Appelbaum PS (1995), \textit{Law and Human Behavior} 19(2):149--174;\\
Health Canada, \textit{Sixth Annual Report on MAiD} (2025);\\
Heijnen WT et al.\ (2010), \textit{J Clin Psychopharmacol} 30(5):616--619;\\
Hindmarch T, Hotopf M, Owen GS (2013), \textit{BMC Medical Ethics} 14:54;\\
Kim SYH, De Vries RG, Peteet JR (2016), \textit{JAMA Psychiatry} 73(4):362--368;\\
Moye J \& Marson DC (2007), \textit{J Gerontol B Psychol Sci Soc Sci} 62(1):P3--P11;\\
Owen GS et al.\ (2008), \textit{BMJ} 337:a448;\\
Owen GS et al.\ (2009), \textit{Br J Psychiatry} 195(3):257--263;\\
Vollmann J et al.\ (2003), \textit{Psychological Medicine} 33(8):1463--1471.
}

\end{document}